\section{Related work}
\label{sec:related}
\paragraph{Multimodal decision models.}
Recent Jev-style systems map runtime-defined alternatives to probabilities. Visual Jev and PixelJev adapt vision--language models using option-token readouts~\citep{yu_2026_visualjev,zhou_2026_pixeljev}, while Koa-action learns atomic label tokens with randomized label mappings~\citep{dai_2026_koa}. OmniJev combines learned option scoring with probabilistic supervision and calibration~\citep{xu_2026_omnijev}, and Chinese-Jev develops a text-encoder decision model with heterogeneous supervision~\citep{wang_2026_chinesejev}. Shared-context execution provides a complementary systems direction: Visual Jev reuses image prefixes, and FastJEV studies state reuse and candidate-layer pruning~\citep{yu_2026_visualjev,ma_2026_fastjev}. OpenVisualJev focuses on the scoring architecture, combining independent candidate representations with explicit, permutation-equivariant joint comparison.

\paragraph{Set-dependent choice.}
Deep Sets establishes invariant and equivariant set representations, and Set Transformer models interactions through attention~\citep{zaheer_2017_deepsets,lee_2019_set}. Context-dependent choice models such as FETA and FATE allow an alternative's utility to depend on the surrounding set, including preference reversals when another alternative changes~\citep{pfannschmidt_2019_context}. Our candidate module applies these principles to multimodal representations. The relation experiments test this capacity directly, alongside the stability of scores under candidate permutations.

\paragraph{Probability quality and decision evaluation.}
Jev benchmarks and decision-gate studies examine calibration, selective prediction, option ordering, and candidate-count robustness~\citep{deusser_2026_benchmark,rafe_2026_gates}. Judge routing, security screening, and mobile action selection illustrate the downstream use of decision scores~\citep{li_2026_jevjudge,liu_2026_security,zhang_2026_jevmobile}. Visual MCQA debiasing and selective VQA likewise connect answer confidence with reliable selection~\citep{atabuzzaman_2025_bias,eisenschlos_2024_selective}. Our evaluation separates recognition accuracy, order stability, and probability quality. The calibration pipeline builds on confidence--accuracy regularization and multiclass calibration objectives~\citep{liang_2020_dca,hebbalaguppe_2022_mdca}, followed by temperature scaling~\citep{guo_2017_calibration}; ablations measure the contributions of training-time and post-training adjustment.

\paragraph{Fine-grained visual recognition.}
FOCI formulates fine-grained recognition as multiple-choice answering with difficult distractors~\citep{geigle_2024_foci}. Finedefics combines attribute-mediated alignment with classification instruction tuning~\citep{he_2025_finedefics}. We reuse the candidate records released with the Finedefics evaluation code and train a unified Qwen3-VL model across six datasets. This provides a common recognition setting for comparing scoring architectures and complements the controlled structural experiments.
